\chapter{전기적 기능에 따른 뉴런의 유형}
\chaptermark{계측기로서의 뉴런}
\label{sec:eetypes}
\begin{chaptergoals}
\item 뉴런이 동작하는 소자들: 적분기, 필터, 변환기, 발진기, 래치;
\item 세포 하나가 대역 통과 공진기가 되거나 \nt{SERO}가 켜는 래치가 되는 방식;
\item 되먹임이 누설 뉴런으로 무누설 적분기를 만드는 방식과 정밀한 조정이 필요한 이유;
\item 세포 유형이 아니라 이온 채널과 브로드캐스트 채널이 정하는 한 세포의 모드인 이유.
\end{chaptergoals}

\ref{sec:neurontypes}장에서 우리는 뉴런을 그 출력이 내보내는 신경전달물질에 따라 분류했다. 전자 공학자라면 다르게 분류할 것이다. 뉴런이 자기 신호에 무엇을 하는지, 즉 어떤 소자로 동작하는지에 따라서다. 이 책의 주요 소자는 \ref{sec:glutcomputer}장부터 알고 있다. 누설 적분 뉴런~\eqref{eq:glutneuron}, 곧 비교기가 달린 RC 회로다. 그러나 뇌에는 다른 소자도 있다. 표~\ref{tab:eetypes}는 이들을 네 무리로 모았으며, 거의 모두 이 책에 이미 나왔다.

\begin{center}
\footnotesize
\setlength{\tabcolsep}{3pt}
\renewcommand{\arraystretch}{1.15}
\begin{longtable}{>{\raggedright}p{3.1cm}>{\raggedright}p{3.3cm}>{\raggedright}p{4.7cm}>{\raggedright\arraybackslash}p{2.4cm}}
\caption{소자로서의 뉴런: 이름, 전자 공학의 대응물, 소자가 하는 일, 책에서 나오는 곳.}\label{tab:eetypes}\\
\toprule
뉴런 & 소자 & 하는 일 & 책에서 나오는 곳 \\
\midrule
\endfirsthead
\toprule
뉴런 & 소자 & 하는 일 & 책에서 나오는 곳 \\
\midrule
\endhead
\bottomrule
\endfoot
\multicolumn{4}{l}{\emph{필터와 적분기}} \\
누설 적분 뉴런 & 비교기가 달린 RC 적분기 & 창 $\tau$ 안의 입력을 더하고 문턱값에서 발화한다 & \ref{sec:glutcomputer}장, \eqref{eq:glutneuron} \\
동시성 검출기 & 시간 창이 있는 AND & 입력이 거의 함께 올 때만 발화한다; 매우 작은 $\tau$ & \ref{sec:glutcomputer}, \ref{sec:code}장, \eqref{eq:window} \\
위상성 뉴런 & 고역 통과 필터, 에지 검출기 & 입력의 변화에 반응하고 곧 잠잠해진다 & \ref{sec:sensors}장 \\
적응 뉴런 & 자동 이득 조절이 달린 고역 통과 필터 & 입력이 일정하면 발화율이 떨어진다 & \ref{sec:sensors}장, \eqref{eq:nakarushton} \\
공진기 & 공진 회로, 대역 통과 필터 & 자기 주파수의 입력에 가장 잘 반응한다 & \ref{sec:resonator}절 \\
무누설 적분기 & 연산 증폭기 적분기 & 속도를 위치로 바꾸고 그 위치를 유지한다 & \ref{sec:perfectint}절 \\
\midrule
\multicolumn{4}{l}{\emph{변환기}} \\
긴장성 뉴런 & 전압---주파수 변환기 & 발화율이 입력을 선형으로 따르고 거의 지치지 않는다 & \ref{sec:code}장 \\
등급 전위 뉴런(스파이크 없음) & 아날로그 증폭기, 버퍼 & 매끄러운 전압을 짧은 거리로 전달한다 & \ref{sec:sensors}장 \\
정류 뉴런 & 반파 정류기 & 발화율은 음수가 될 수 없다, $[u]_+$ & \ref{sec:glutcomputer}, \ref{sec:gaba}장 \\
“켜짐---꺼짐” 쌍 & 정류기 쌍, 이중 레일 부호 & 하나는 “더 크다”에, 다른 하나는 “더 작다”에 반응한다 & \ref{sec:glutcomputer}장, \eqref{eq:dualrail} \\
분로 억제 뉴런 & 아날로그 나눗셈기, 이득 조절 & 입력에서 빼지 않고 입력을 나눈다 & \ref{sec:gaba}장, \eqref{eq:shunt} \\
\midrule
\multicolumn{4}{l}{\emph{발진기와 시간}} \\
페이스메이커 & 이완 발진기, 멀티바이브레이터 & 스스로 일정한 주파수로 스파이크를 낸다 & \ref{sec:serocomputer}, \ref{sec:dofa}장 \\
입력이 있는 페이스메이커 & 전압 제어 발진기 & 자기 리듬이 있고 입력이 이를 빠르게 하거나 늦춘다 & \ref{sec:chemoutput}장 \\
버스트 뉴런 & 게이트형 버스트 발생기 & 잦은 스파이크의 버스트, 즉 “대시”를 낸다 & \ref{sec:code}장, \eqref{eq:burst} \\
위상 고정 뉴런 & 위상 검출기, 위상 고정 루프 & 입력 파동의 특정 위상에서 발화한다 & \ref{sec:sensors}, \ref{sec:chemoutput}장 \\
지연 축삭 & 지연선 & 신호를 정해진 시간만큼 늦춘다 & \ref{sec:glutcomputer}장, 그림~\ref{fig:itd} \\
\midrule
\multicolumn{4}{l}{\emph{기억과 전환}} \\
쌍안정 뉴런 & 슈미트 트리거, 래치 & 짧은 입력이 오래가는 상태로 전환한다 & \ref{sec:bistable}절 \\
가지돌기 가지가 있는 뉴런 & 멀티플렉서, 작은 다층 네트워크 & 허용 입력이 있을 때만 가지가 입력을 통과시킨다 & \ref{sec:synthesis}장 \\
\end{longtable}
\end{center}

표 전체보다 중요한 단서가 하나 있다. 이것들은 서로 다른 세포가 아니라 서로 다른 \emph{모드}다. 같은 뉴런이 느린 입력에서는 적분기가 되고, 억제가 $\tau$를 짧게 하면 동시성 검출기가 된다(\ref{sec:code}장). 깨어 있을 때는 페이스메이커이고, 잘 때는 조용한 수신기다. 어떤 소자가 될지는 막의 이온 채널과, 그 채널을 조정하는 브로드캐스트 채널이 정한다.

\section{이 책에 아직 나오지 않은 네 소자}

\subsection{공진기}
\label{sec:resonator}

어떤 뉴런에는 전압의 모든 변화에 맞서는 느린 전류가 있다. 전압이 오르면 끌어내리고 내리면 끌어올리지만, 지연이 있다. 전자 공학자에게 이런 전류는 인덕턴스처럼 행동하며, 막의 커패시턴스와 함께 공진 회로를 이룬다. 뉴런은 특정 주파수, 보통 몇 Hz에서 수십 Hz의 입력에 가장 강하게 반응하고, 그보다 느리거나 빠른 입력에는 약하게 반응한다~\cite{HutcheonYarom2000}. 세포 하나로 된 대역 통과 필터다. 잡음 섞인 입력에서 자기 주파수의 리듬, 예를 들어 \ref{sec:gaba}장의 국소 리듬을 골라낸다.

\subsection{무누설 적분기}
\label{sec:perfectint}

누설 적분기는 입력을 시간 $\tau$, 즉 수십 밀리초 동안만 기억한다. 그러나 눈이 움직이지 않고 머물려면 자기 위치를 몇 초 동안 기억해야 한다. “돌려라”라는 명령은 짧은 속도로 오지만, 눈은 새 위치에 붙들어 두어야 한다. 네트워크는 이 문제를 \ref{sec:glutcomputer}장의 기억과 같은 되먹임으로 푼다. 시상수 $\tau$인 뉴런 무리가 가중치 $w$로 자신을 흥분시키면 $\tau\,\dd r/\dd t=-r+w\,r+I$이고, 무리의 시상수는
\begin{empheq}[box=\keybox]{equation}
\tau_{\text{eff}} \;=\; \frac{\tau}{1-w}
\label{eq:perfectint}
\end{empheq}
이다. $\tau=0.1\,\text{s}$, $w=0.995$이면 $20\,\text{s}$가 된다. 하나하나는 10분의 1초 만에 잊는 뉴런들로 만든 거의 이상적인 적분기다~\cite{Seung1996}. 대가는 정밀한 조정이다. $w$가 1보다 조금만 커도 적분기는 포화로 치닫고, 조금만 작아도 빨리 잊는다. 그래서 이런 적분기에는 \ref{sec:motorlearning}장에서 설명한 것 같은 학습에 의한 끊임없는 조정이 필요하다.

\subsection{쌍안정 뉴런}
\label{sec:bistable}

\ref{sec:glutcomputer}장과 \ref{sec:gaba}장에서는 플립플롭을 뉴런 무리로 만들었다. 그러나 세포 하나로 된 플립플롭도 있다. 어떤 뉴런에는 한번 켜지면 스스로 흥분을 유지하는 전류가 있다. 짧은 입력이 뉴런을 오래가는 발화 상태, 즉 \emph{플래토}로 옮기고, 억제가 다시 휴지 상태로 되돌린다. 이것은 슈미트 트리거다. 뉴런에는 안정한 상태가 둘 있고 그 사이에 히스테리시스가 있다. 예를 들어 근육을 제어하는 \nt{ACET} 뉴런이 이렇게 행동하며, 이 성질은 브로드캐스트 채널이 켠다. 세로토닌이 뉴런에 닿으면 플래토가 나타난다~\cite{Hounsgaard1988}. 같은 세포가 \nt{SERO}가 있으면 래치이고, 없으면 평범한 적분기다.

\subsection{적분기와 공진기}

이론은 흥분성 세포를 두 부류로 나눈다~\cite{Izhikevich2007}. \emph{적분기}는 0에서 시작하는 전압 제어 발진기를 닮았다. 문턱값을 조금 넘으면 뉴런은 얼마든지 느리게 발화할 수 있고, 발화율은 입력에 따라 매끄럽게 커진다. 이런 뉴런은 들어온 것을 모두 더하며, 입력의 크기를 발화율로 잘 전달한다. \emph{공진기}는 최소 주파수가 있는 발진기를 닮았다. 느리게 발화하지 못하고, 문턱값 수준의 입력에서 곧바로 유한한 주파수로 스파이크를 내며, 자기 주파수의 입력을 선호한다. 앞의 것은 \ref{sec:code}장의 발화율 부호에 자연스러운 소자이고, 뒤의 것은 타이밍 부호에 자연스러운 소자다.

\begin{keyblock}
\begin{proposition}[뉴런 하나, 여러 소자]
\label{prop:eetypes}
전기적 기능으로 보면 뉴런은 누설 적분기와 무누설 적분기, 동시성 검출기, 고역 통과 필터와 대역 통과 필터, 전압---주파수 변환기, 정류기, 나눗셈기, 발진기, 지연선, 플립플롭으로 동작한다. 주어진 세포가 어떤 소자가 될지는 그 세포의 이온 채널과, 그 채널을 조정하는 브로드캐스트 채널이 정한다. 모드 자체는 측정되었다. 뇌가 이런 재구성을 계산에 얼마나 쓰는지는 대부분 가설이다.
\end{proposition}
\end{keyblock}

엔지니어에게 이것은 프로그래머블 로직 어레이와 비슷하다. 하드웨어는 하나이고 소자는 구성이 정한다. 다만 여기서는 구성을 펌웨어를 쓸 때가 아니라 동작 중에 바꾼다. \ref{sec:serocomputer}--\ref{sec:acet}장에서 다룬 느린 브로드캐스트 채널이 바꾼다.

\section{연습문제}

\begin{exercise}[\lvlA]
\label{ex:18:1}
\emph{무누설 적분기의 허용 오차.} $\tau=0.1$~s인 뉴런들이 식~\eqref{eq:perfectint}에 따라 눈의 위치를 $T=1$~s 동안 양쪽으로 $5\%$ 이내의 오차로 유지한다. (a)~허용되는 $w$의 범위를 구하라. (b)~$w$는 시냅스 $K$개의 가중치 합이고 각각 독립적인 $10\%$ 오차가 있다. $K$가 얼마일 때 합의 편차가 허용 범위에 들어가는가?
\end{exercise}

\begin{exercise}[\lvlB]
\label{ex:18:2}
\emph{회로로서의 공진기.} \ref{sec:resonator}절의 느린 전류를 변수 $W$로 나타내자: $C\,\dd V/\dd t=-g_LV-g_wW+I$, $\tau_w\,\dd W/\dd t=V-W$. 이것이 $R_w=1/g_w$, $L_w=\tau_w/g_w$인 병렬 가지를 가진 $RC$ 회로임을 보이고, $C/g_L=10\ms$, $\tau_w=100\ms$, $\alpha=g_w/g_L=4$에서 공진 주파수와 그곳에서 $|Z|$가 $|Z(0)|$보다 몇 배 큰지 구하라.
\end{exercise}

\begin{exercise}[\lvlB]
\label{ex:18:3}
\emph{적분기는 어떻게 발화를 시작하는가.} (a)~뉴런 $\tau\,\dd V/\dd t=-V+\mu$, $\tau=10\ms$, 문턱값 $\theta$, 0으로 재설정. $\mu/\theta-1$이 얼마일 때 발화율이 $1$~Hz인가? (b)~모델 $\tau\,\dd v/\dd t=v^2+\mu$(스파이크는 $v$가 $+\infty$로 가는 것, 재설정은 $-\infty$로)는 $\mu=0.01$에서 발화율이 얼마인가?
\end{exercise}

\begin{exercise}[\lvlB]
\label{ex:18:4}
\emph{모델링: 적분기와 공진기.} $g_L=1$, $\tau_m=10\ms$, $\tau_w=100\ms$, 문턱값 $1$, 재설정 $V\to0$($W$는 그대로)인 연습문제~\ref{ex:18:2}의 모델에 $\mu=1.5\sin2\pi ft$, $f=1$--$40$~Hz를 넣는다. $\alpha=0$과 $\alpha=4$에서 뉴런은 어떤 주파수 대역에서 스파이크를 내는가? 조건 $1.5\,|Z(f)|>1$과 비교하라.
\end{exercise}

\begin{exercise}[\lvlB]
\label{ex:18:5}
\emph{세포 하나로 된 래치.} 플래토 전류가 있는 뉴런: $\tau\,\dd r/\dd t=-r+F(wr+I)$, $F(x)=\min(1,\max(0,x))$, $\tau=10\ms$, $w=1.5$, 배경 $I=-0.2$. (a)~$I=0.3$ 펄스가 뉴런을 $r=0$에서 플래토로 옮기는 최소 길이는? (b)~긴 펄스 $I=-0.2-B$가 뉴런을 휴지 상태로 되돌리는 최소 진폭 $B$는?
\end{exercise}

\begin{exercise}[\lvlB]
\label{ex:18:6}
\emph{적분기의 자기 조정.} 시선을 고정할 때 입력은 $I=0$이고, 미끄러짐 센서가 표류 속도 $e=\dd r/\dd t$를 주며, $\dd w/\dd t=-\eta\,e\,r$이다. $w=0.95$, $\tau=0.1$~s, $\langle r^2\rangle=1$에서 고정 $60$~s 뒤 $|1-w|$가 연습문제~\ref{ex:18:1}의 허용 범위에 드는 $\eta$를 구하라. 눈을 돌리는 동안에도 학습하면 무엇이 망가지는가?
\end{exercise}

\begin{exercise}[\lvlC]
\label{ex:18:7}
\emph{소자를 왜 재구성하는가?} 브로드캐스트 채널이 연습문제~\ref{ex:18:2} 모델의 $\alpha$를 $0$과 $4$ 사이에서 바꾼다. 전류의 백색 잡음 속 약한 사인파에 대해, $11$~Hz와 $1$~Hz에서 공진기는 적분기보다 신호 대 잡음비가 몇 배 좋은가? 바로 모드 전환이 유리한 과제를 고안하라.
\end{exercise}
